\subsection{The cap-free harmonic core}
\label{sec:lower}

This subsection develops the cap-free construction used by the last two
branches of the revealing curve.  Its notation is local to this subsection.
We first record its extremal specialization and then extract the scaled form
reused by the mixed branch.  The construction
has many equally large groups of positive jobs, revealed from longest to
shortest, followed by one large group of zero jobs.  The only design choice is
the gap between two consecutive positive processing times.  Choosing these
gaps harmonically will make all decisions of the online algorithm disappear
from the leading term.

\begin{theorem}[Extremal harmonic core]
\label{thm:lower}
For every deterministic online algorithm $\mathcal B$ and every $\eps>0$,
there are arbitrarily large fixed instances $I$ with unit tests and rational
processing times such that
\[
  \ALG_{\mathcal B}(I)\ge (\RdetOT-\eps)\OPT(I).
\]
\end{theorem}

\paragraph{The revelation experiment.}
Fix an integer $K$, a rational number $\xi>0$, and a rational slack
$\gamma\in(0,1)$; one may keep $\gamma=1/2$ throughout.  For a large scaling
integer $H$, the instance will contain $H$ jobs of each of $K$ positive types
$0,1,\ldots,K-1$ and $\xi H$ zero jobs.  We choose $H$ to clear all
denominators.

The adversary assigns processing times according to the order in which tests
are started.  The first $H$ tested jobs receive processing time $p_0$, the
next $H$ receive $p_1$, and so on; after the $K$ positive blocks, every
remaining job receives processing time zero.  Thus
\[
  p_0>p_1>\cdots>p_{K-1}>1,
\]
so the algorithm sees the positive types from longest to shortest and learns
about the zero jobs only at the end.

This adaptive description is only a convenient way to construct the input.
After running the interaction against a fixed deterministic algorithm, assign
to every job label the value revealed when that label was tested.  Replaying
the algorithm on this fixed assignment produces exactly the same transcript.
If the algorithm never tests or completes some job, its cost is infinite and
the lower bound is immediate; otherwise every job label is recorded.
Hence it is enough to lower-bound the cost in the revelation experiment.

\paragraph{Why each phase can be relaxed to shortest processing time.}
Let $T_\ell$ be the start of the first test in positive block $\ell$, and let
$T_K$ be the start of the first test in the zero block.  Phase $\ell$ contains
the completions in $(T_\ell,T_{\ell+1}]$; a completion exactly at $T_\ell$
belongs to the preceding phase.  Each $T_\ell$ is the boundary between two
nonpreemptive operations.  Consequently, a previously tested job that is
unfinished at $T_\ell$ has not started processing and still needs its entire
processing time.

We repeatedly use the following elementary observation.  Suppose jobs with
remaining work $v_1,\ldots,v_m$ complete after some time $T$.  If their
completion order is $q_1,\ldots,q_m$, then the $r$th relative completion time
is at least $\sum_{h\le r}v_{q_h}$.  Summing these prefix inequalities shows
that their total relative completion time is minimized by arranging the
$v_i$ in nondecreasing order.  This is merely a lower bound; the online
algorithm need not actually use that order.

For normalized group masses $q_1,\ldots,q_m$ and remaining lengths
$v_1\le\cdots\le v_m$, this shortest-first bound has leading coefficient
\begin{equation}
 \mathcal S((q_g,v_g)_{g=1}^m)
 =\sum_{g=1}^m
   \left(\frac{v_gq_g^2}{2}
    +q_g\sum_{h<g}v_hq_h\right).
 \label{eq:lower-S}
\end{equation}
More precisely, the corresponding contribution for groups of size $q_gH$
is $H^2\mathcal S+O(H)$.  The quadratic term is simply the familiar sum
$v_g(1+\cdots+q_gH)$ inside one group; the cross term is the waiting caused
by all shorter groups.

We ensure that
\begin{equation}
  p_0-p_{K-1}<1.
  \label{eq:lower-span}
\end{equation}
In phase $\ell$, an unfinished job of an older type $j<\ell$ therefore has
remaining length $p_j$, whereas a type-$\ell$ job completed in its own phase
needs both its test and processing, of total length $1+p_\ell$.  Their order
in the shortest-first relaxation is
\[
  p_{\ell-1}<p_{\ell-2}<\cdots<p_0<1+p_\ell.
\]
At $T_K$, an untested zero job has remaining length $1$, so the relaxed tail
starts with all zero jobs and then uses the positive types from shortest to
longest.  These are the only two places where
\cref{eq:lower-span} is used.

\paragraph{The finite accounting identity.}
We now record the online algorithm's possible behavior.  Normalize all job
counts by $H$.  Let
\begin{itemize}
\item $s_j\in[0,1]$ be the mass of type-$j$ jobs completed in their own
      phase;
\item $d_{j,\ell}\ge0$ be the mass of type-$j$ jobs completed in the later
      phase $\ell>j$; and
\item
\[
  X_{j,\ell}=s_j+\sum_{r=j+1}^{\ell-1}d_{j,r}\le1
  \qquad(\ell>j)
\]
be the mass of type $j$ completed before phase $\ell$.  Thus
$X_{j,\ell+1}$ is the mass completed by the end of phase $\ell$.
\end{itemize}
At time $T_\ell$, the machine has already performed the tests of the $\ell H$
jobs in earlier positive blocks and has processed every earlier job that has
completed.  Thus
\begin{equation}
  \frac{T_\ell}{H}
  \ge \ell+\sum_{j<\ell}p_jX_{j,\ell}.
  \label{eq:lower-start}
\end{equation}
For each phase, charge its completions their common start time from
\cref{eq:lower-start}, and charge their relative completion times using
\cref{eq:lower-S}.  Do the same at $T_K$ for the remaining positive jobs
and the $\xi H$ zero jobs.  Expanding and collecting equal variables gives
\begin{align}
 \frac{\ALG}{H^2}
 \ge{}& C_K
 +\sum_{j=0}^{K-1}\left(\frac{s_j^2}{2}-\theta_js_j\right)
 \nonumber\\
 &+\sum_{j<\ell}d_{j,\ell}
 \left[
   \ell-j-\theta_j
   -\sum_{j<k<\ell}(p_j-p_k)X_{k,\ell+1}
 \right]-O(1/H),
 \label{eq:lower-accounting}
\end{align}
where $K,\xi,\gamma$ are fixed while $H\to\infty$, and
\begin{align}
 \theta_j
   &=1-\xi(p_j-1)-\sum_{k>j}(p_j-p_k-1),
   \label{eq:lower-theta}\\
 C_K
   &=\frac{\xi^2}{2}+2\xi K+K^2+P_K,
 &
 P_K&=\sum_{j=0}^{K-1}\left(j+\frac12\right)p_j.
 \label{eq:lower-CP}
\end{align}

For completeness, here is how to audit
\cref{eq:lower-accounting} without guessing its origin.  In phase $\ell$,
the masses in the shortest-first relaxation are
\[
 d_{\ell-1,\ell},d_{\ell-2,\ell},\ldots,d_{0,\ell},s_\ell
\]
with respective lengths
\[
 p_{\ell-1},p_{\ell-2},\ldots,p_0,1+p_\ell.
\]
Multiply their total mass by the right-hand side of
\cref{eq:lower-start} and add the functional
\cref{eq:lower-S}.  In the final phase use masses
\[
 \xi,\ 1-X_{K-1,K},\ldots,1-X_{0,K}
\]
and lengths $1,p_{K-1},\ldots,p_0$.  The constant terms sum to $C_K$.
The terms containing only $s_j$ are $s_j^2/2-\theta_js_j$; the coefficient
of each $d_{j,\ell}$ is precisely the bracket in
\cref{eq:lower-accounting}.  Thus the displayed identity is just the
phasewise shortest-first bound with the phase start times retained.  No
optimality assumption about the online schedule is hidden in it.

The meaning of the expression is useful.  Completing an $s_j$ fraction of a
type in its own phase creates a quadratic completion cost but saves later
waiting cost, producing $s_j^2/2-\theta_js_j$.  A completion postponed to a
later positive phase is represented by $d_{j,\ell}$ and its bracket.  We next
choose the processing times so that every such bracket is nonnegative and all
$\theta_j$ are equal.

\paragraph{The harmonic choice and its telescoping cancellation.}
Set
\begin{equation}
  p_{K-1}=1+\frac{\gamma}{\xi},
  \qquad
  p_i=p_{i+1}+\frac1{\xi+K-1-i}
  \quad(0\le i<K-1).
  \label{eq:lower-harmonic}
\end{equation}
In other words, measured upward from the shortest type, the consecutive gaps
are
\[
  \frac1{\xi+1},\frac1{\xi+2},\ldots,\frac1{\xi+K-1}.
\]
This is the key design choice.  If $m=K-1-j$, then
\begin{align*}
 p_j-1
   &=\frac\gamma \xi+\sum_{r=1}^m\frac1{\xi+r},\\
 \sum_{k>j}(p_j-p_k)
   &=\sum_{r=1}^m\frac{r}{\xi+r}.
\end{align*}
Adding the first line multiplied by $\xi$ to the second line makes each
denominator cancel:
\begin{equation}
 \xi(p_j-1)+\sum_{k>j}(p_j-p_k)
 =\gamma+\sum_{r=1}^m\frac{\xi+r}{\xi+r}
 =\gamma+m.
 \label{eq:lower-telescope}
\end{equation}
Substitution in \cref{eq:lower-theta} gives the same value for every type,
\begin{equation}
  \theta_j=1-\gamma.
  \label{eq:lower-theta-constant}
\end{equation}
This is why the increments in \cref{eq:lower-harmonic} are harmonic: the
factor $\xi$ coming from the zero block and the multiplicity $r$ coming from the
$r$ shorter positive types add up to the denominator $\xi+r$.

The delayed-completion terms are now harmless as well.  Since
$X_{k,\ell+1}\le1$, their brackets are at least
\begin{align*}
 &\ell-j-(1-\gamma)
   -\sum_{j<k<\ell}(p_j-p_k)\\
 &\hspace{20mm}
 =\gamma+\sum_{j<k<\ell}\bigl(1-(p_j-p_k)\bigr)>0,
\end{align*}
where the final inequality follows from
\cref{eq:lower-span}.  We may therefore drop every term containing
$d_{j,\ell}$.  The remaining minimization is not a multivariable optimization
problem: complete one square for each $j$,
\[
  \frac{s_j^2}{2}-(1-\gamma)s_j
  =\frac12\bigl(s_j-(1-\gamma)\bigr)^2
   -\frac12(1-\gamma)^2.
\]
Because $1-\gamma\in(0,1)$, this minimum is feasible for
$s_j\in[0,1]$.  Hence every online algorithm obeys
\begin{equation}
 \liminf_{H\to\infty}\frac{\ALG}{H^2}
 \ge Q_K
 \defeq C_K-\frac K2(1-\gamma)^2.
 \label{eq:lower-QK}
\end{equation}

It remains to compute the offline benchmark.  Knowing all processing times,
the optimum runs the zero jobs first and the positive types from shortest to
longest.  Equivalently, the pairwise contribution of two jobs is
$1+\min\{p_i,p_j\}$.  Its leading coefficient is
\begin{equation}
  D_K=\frac{\xi^2}{2}+\xi K+\frac{K^2}{2}+P_K,
  \qquad
  \lim_{H\to\infty}\frac{\OPT}{H^2}=D_K.
  \label{eq:lower-DK}
\end{equation}
The four terms respectively account for zero--zero pairs, zero--positive
pairs, the unit-test part of positive--positive pairs, and the processing
part of positive--positive pairs.  Thus the finite-$K$ lower-bound
certificate is simply $Q_K/D_K$.

\paragraph{Passing to a smooth curve.}
Write
\[
  \alpha=\frac K\xi
\]
for the ratio of positive mass to zero mass, and keep $\alpha$ fixed while
$K\to\infty$.  We restrict to $0<\alpha<e-1$.  Then
\[
 p_0-p_{K-1}
 =\sum_{r=1}^{K-1}\frac1{\xi+r}
 <\ln\!\left(1+\frac K\xi\right)
 =\ln(1+\alpha)<1,
\]
which verifies \cref{eq:lower-span}.

The processing contribution $P_K$ has the exact form
\begin{equation}
 P_K
 =\frac{K^2}{2}\left(1+\frac\gamma \xi\right)
  +\frac12\sum_{h=1}^{K-1}\frac{(K-h)^2}{\xi+h}.
 \label{eq:lower-P-exact}
\end{equation}
Indeed, the common baseline $1+\gamma/\xi$ contributes the first term, and
swapping the order of summation shows that the increment $1/(\xi+h)$ receives
total weight $(K-h)^2/2$.  After division by $\xi^2$, the sum becomes a
Riemann integral:
\begin{align}
 \frac{P_K}{\xi^2}&\longrightarrow I(\alpha)
   =\frac{\alpha^2}{2}
    +\frac12\int_0^\alpha\frac{(\alpha-t)^2}{1+t}\,dt
   \nonumber\\
 &=\int_0^\alpha(\alpha-t)\bigl(1+\ln(1+t)\bigr)\,dt
   \nonumber\\
 &=\frac12(1+\alpha)^2\ln(1+\alpha)
   -\frac\alpha2-\frac{\alpha^2}{4}.
 \label{eq:lower-I}
\end{align}
The second integral form follows either by integration by parts or by summing
the limiting processing-time curve $1+\ln(1+t)$ against the remaining
positive mass $\alpha-t$.
The subtraction in \cref{eq:lower-QK} is only $O(K)$, whereas
$C_K,D_K=\Theta(K^2)$.  Since $C_K-D_K=\xi K+K^2/2$, we obtain
\begin{equation}
 R(\alpha)
 \defeq\lim_{K\to\infty}\frac{Q_K}{D_K}
 =1+
 \frac{\alpha+\alpha^2/2}
 {\frac12+\frac\alpha2+\frac{\alpha^2}{4}
  +\frac12(1+\alpha)^2\ln(1+\alpha)}.
 \label{eq:lower-Ralpha}
\end{equation}

This last one-variable maximization is where the constant in the theorem
comes from.  Put $z=(1+\alpha)^2$.  Then
\begin{equation}
  R(z)=1+\frac{2(z-1)}{1+z+z\ln z}.
  \label{eq:lower-Rz}
\end{equation}
The derivative of the fraction has the sign of
\[
  3-z+\ln z.
\]
For $z>1$ this expression is strictly decreasing, is positive at $z=1$, and
tends to $-\infty$.  Hence it has one zero $\zstar$, characterized by
\[
  \zstar-3=\ln \zstar.
\]
The derivative is positive before $\zstar$ and negative afterward, so this is
the global maximum.  At the root,
\[
  1+\zstar+\zstar\ln \zstar=(\zstar-1)^2,
\]
and therefore
\[
  R(\zstar)=1+\frac{2}{\zstar-1}=\RdetOT.
\]
Moreover, $\sqrt{\zstar}-1=1.12255\ldots<e-1$, so the maximizer lies in the
range where \cref{eq:lower-span} holds.

\begin{proof}[Proof of \cref{thm:lower}]
Choose a rational $\alpha$ sufficiently close to
$\alpha_\star=\sqrt{\zstar}-1$, and choose any rational
$\gamma\in(0,1)$.  For a sufficiently large $K$, set $\xi=K/\alpha$ and use
\cref{eq:lower-harmonic}; then $Q_K/D_K>\RdetOT-\eps/2$.
All parameters are rational.  Let $H$ tend to infinity through multiples of
their common denominators.  \Cref{eq:lower-QK,eq:lower-DK} make the actual ratio exceed
$\RdetOT-\eps$ for all sufficiently large $H$.  Finally, the replay argument
at the beginning of the proof converts the revelation experiment into a fixed
instance for the given deterministic algorithm.
\end{proof}

\begin{lemma}[Scaled harmonic core]
\label{lem:opt:harmonic-block}
Fix a total normalized mass $s>0$ and a parameter $t\in[1,e)$.  Against
every deterministic obligatory-testing algorithm, there are scales
$M_\nu\to\infty$ and finite rational multilevel instances with
$sM_\nu+o(M_\nu)$ jobs such that
\begin{equation}
 \frac{\OPT}{M_\nu^2}\longrightarrow
 O_0=\frac{s^2}{4}(1+t^{-2}+2\ln t),
 \qquad
 \liminf_{\nu\to\infty}\frac{\ALG}{M_\nu^2}
 \ge A_0^{\rm lb}
 \defeq O_0+\frac{s^2}{2}(1-t^{-2}).
 \label{eq:opt:harmonic-coefficients}
\end{equation}
Their normalized total positive processing mass converges to
\begin{equation}
 \frac1{M_\nu}\sum_{j:p_j>0}p_j\longrightarrow L=s\ln t,
 \label{eq:opt:harmonic-processing-mass}
\end{equation}
and their largest processing time converges to $1+\ln t<2$.
The adaptive revelation is converted by replay into a fixed instance.
\end{lemma}

\begin{proof}
The case $t=1$ is obtained by continuity from instances with vanishing
positive mass, so suppose $1<t<e$.  In the construction above set
$\alpha=t-1=K/\xi$.  Before rescaling, the total normalized mass is
$K+\xi=\xi t$.  Give every unit of normalized mass the common factor
$s/(\xi t)$.

The phase accounting, harmonic cancellation, and completion of squares in
\cref{eq:lower-accounting,eq:lower-telescope,eq:lower-QK} are homogeneous of
degree two in the job masses.  The Riemann-sum limit
\cref{eq:lower-I} substituted into \cref{eq:lower-DK} therefore gives
\[
 O_0=\frac{s^2}{4}(1+t^{-2}+2\ln t).
\]
Likewise $C_K-D_K=\xi K+K^2/2$, while the $O(K)$ square-completion loss in
\cref{eq:lower-QK} disappears after division by the quadratic scale.  After
the same rescaling this gives the guaranteed lower benchmark
\[
 A_0^{\rm lb}-O_0=\frac{s^2}{2}(1-t^{-2}).
\]
Finally, summing the limiting processing-time curve
$1+\ln(1+x)$ over the positive mass gives $L=s\ln t$, and
\cref{eq:lower-harmonic} gives $p_0\to1+\ln t$.  Rational approximation,
clearing denominators, and the replay argument used in the theorem produce a
diagonal sequence of arbitrarily large fixed rational instances with the
stated limit and liminf bounds.
\end{proof}
